% Заключение: кратко фиксируем достигнутые результаты и перспективы.
Проведенное исследование подтвердило, что задача персонализации образовательной траектории может быть решена в рамках воспроизводимого программного контура, если методическая часть и инженерная реализация проектируются согласованно. В теоретической части работы обоснована актуальность адаптивного обучения в условиях цифровой трансформации образования и показано, что накопление учебных данных открывает возможность перехода от усредненного сценария к персонализированному управлению учебным процессом. В методическом блоке сформирована модель, объединяющая вероятностное отслеживание знаний, алгоритм последовательного выбора педагогического действия и модуль объяснения рекомендаций.

Практическая реализация продемонстрировала применимость выбранного подхода в типовом сценарии цифровой образовательной платформы. Была реализована модульная архитектура, включающая контур сбора событий, аналитический слой, рекомендательный модуль и интерфейсы для разных ролей пользователей. Важным результатом стало формирование управляемого цикла "данные -- решение -- обратная связь", где преподаватель и эксперт сохраняют возможность корректировать работу алгоритма. Таким образом, система решает не только задачу автоматизации рекомендаций, но и задачу прозрачной интеграции алгоритмического слоя в педагогическую практику.

По итогам демонстрационных запусков подтверждена положительная динамика ключевых метрик качества рекомендаций, а также практическая устойчивость сценариев взаимодействия для студентов, методистов и экспертов. Анализ логов показал, что контур экспертной корректировки повышает интерпретируемость решений и позволяет быстрее устранять локальные ошибки политики. Это особенно значимо для образовательных систем, где цена некорректной рекомендации выражается не только в снижении формальной метрики, но и в потере мотивации обучающегося.

Полученные результаты позволяют считать цель исследования достигнутой: разработан и описан прототип адаптивной системы обучения, использующий методы машинного обучения для персонализации контента и темпа обучения. Дополнительно сформирован практический шаблон, пригодный для адаптации под новые дисциплины за счет конфигурационного подхода и модульного разделения ответственности компонентов.

Дальнейшее развитие проекта целесообразно вести по нескольким направлениям. Во-первых, необходимо расширить эмпирическую базу и провести длительное тестирование на реальных учебных потоках с контролем стабильности результатов в течение семестра. Во-вторых, перспективно усилить модель объяснения рекомендаций, включая более строгую валидацию текстовых обоснований и учет педагогического контекста дисциплины. В-третьих, важной задачей остается интеграция с внешними LMS и корпоративными образовательными платформами через стандартизированные интерфейсы обмена событиями. Комплексное развитие этих направлений позволит перейти от демонстрационного прототипа к промышленной эксплуатации адаптивной системы в университетской и профессиональной образовательной среде.
Также перспективно внедрение сценариев междисциплинарной персонализации с единым профилем компетенций студента.
